\documentclass[paper=a4,fontsize=10pt]{scrartcl}

% ---------------------------
% Formato general
% ---------------------------
\usepackage[top=2.5cm,bottom=2.5cm,left=2.5cm,right=2.5cm]{geometry}
\usepackage[spanish]{babel}
\usepackage{setspace}
\usepackage[protrusion=true,expansion=true]{microtype}
\usepackage{amsmath,amsfonts,amsthm}
\usepackage{booktabs}
\usepackage{float}
\usepackage{xcolor}
\usepackage{caption}
\usepackage{fancyhdr}
\usepackage{lastpage}
\usepackage{pgfplots}
\pgfplotsset{compat=1.18}

% ---------------------------
% Encabezado y pie
% ---------------------------
\pagestyle{fancy}
\lhead{}
\chead{}
\rhead{}
\lfoot{}
\cfoot{}
\rfoot{\footnotesize Pagina \thepage\ de \pageref{LastPage}}
\renewcommand{\headrulewidth}{0pt}
\renewcommand{\footrulewidth}{0.4pt}

% ---------------------------
% Titulos (tamano moderado)
% ---------------------------
\usepackage{sectsty}
\sectionfont{\normalsize\bfseries}
\subsectionfont{\small\bfseries}

\title{\textbf{Memoria de Calculo Estructural}}
\author{[Nombre del proyecto]}
\date{\today}

\begin{document}
\onehalfspacing
\maketitle
\thispagestyle{fancy}

\section*{Resumen}
Describir brevemente el alcance del analisis y las verificaciones principales.

\section*{1. Datos generales}
\begin{itemize}
  \item Proyecto: [Nombre]
  \item Ubicacion: [Ciudad]
  \item Norma: [E.030 u otra]
  \item Software: [ETABS / SAP2000 / otro]
\end{itemize}

\section*{2. Parametros de analisis}
\begin{table}[H]
\centering
\caption{Parametros principales}
\begin{tabular}{lc}
\toprule
Parametro & Valor \\
\midrule
$Z$ & [0.35] \\
$U$ & [1.0] \\
$S$ & [1.15] \\
$R$ & [5.4] \\
$T_p$ & [0.6 s] \\
$T_l$ & [2.0 s] \\
\bottomrule
\end{tabular}
\end{table}

\section*{3. Formulacion}
\[
S_a = \frac{Z \cdot U \cdot C \cdot S}{R}\,g
\]

\section*{4. Espectro sismico}
\begin{figure}[H]
\centering
\begin{tikzpicture}
\begin{axis}[
width=0.9\linewidth,
height=6.5cm,
xlabel={$T$ (s)},
ylabel={$S_a/g$},
xmin=0, xmax=2.5,
ymin=0, ymax=0.25,
grid=major
]
% Reemplazar por funciones/valores reales del proyecto
\addplot[blue,domain=0:0.6,samples=2] {0.18634};
\addplot[blue,domain=0.6:2.0,samples=120] {0.111804/x};
\addplot[blue,domain=2.0:2.5,samples=120] {0.223608/(x^2)};
\end{axis}
\end{tikzpicture}
\caption{Espectro de diseno}
\end{figure}

\section*{5. Resultados y verificaciones}
\begin{table}[H]
\centering
\caption{Verificacion de cortante basal}
\begin{tabular}{lc}
\toprule
Concepto & Valor \\
\midrule
$V_{\text{est}}$ & [429.36 ton] \\
$V_{\text{din}}$ & [249.95 ton] \\
Factor de escala $\gamma$ & [1.546] \\
\bottomrule
\end{tabular}
\end{table}

\section*{6. Conclusiones}
Indicar cumplimiento normativo, controles que gobiernan el diseno y recomendaciones.

\begin{thebibliography}{9}
\bibitem{e030}
Reglamento Nacional de Edificaciones. Norma E.030: Diseno Sismorresistente.
\end{thebibliography}

\end{document}

